\cotefile{23}
\scanpage{33}
$(u,u')$ est une connexion de type $\unsure{\Lambda|1}$ de $X$ \unsure{en|$-$|$\to$} $X'$. \unsure{Cela|Ceci|On} dit, on aura pour $g\in\Gamma(X/Y)$
\[
 h^{\{Y_{S'}\}}\bigl(\widetilde w(g)\bigr)=\widetilde w'(h\circ g)
\]

3°) Compatibilité \unsure{au produit|aux produits|un produit} : \deleted{\unsure{connexion|idem}}\inserted{soient} $w=(u,v)$, $w'=(u',v)$ deux connexions sur $X/Y$, $X'/Y$, (avec \unsure{$h=v$|$v=v$}), on en déduit une connexion sur $X\times_Y X'$, \deleted{\unsure{un}} i.e. une $S'$-section \unsure{couplée}, \unsure{une|sur} de $(X_{\deleted{\unsure{S'}}}\times_Y X')_{S'}=X_{S'}\times_{Y_{S'}}X'_{S'}$ au dessus de $Y_{S'}$, induisant $v_{\deleted{\unsure{S'}}}$ sur $Y_{S'}$ (\unsure{Morphisme} de structure) \inserted{soit $w''$,} et on aura alors
\margin{\diagram{\begin{tikzcd}[ampersand replacement=\&,column sep=tiny,row sep=tiny]
X_{S'},\arrow[rr,bend right=22] \& X'_{S'}\arrow[rr,bend right=22] \& X_{\unsure{S'|s'}}, \& X'_{S'}\\
Y_{S'}\arrow[rr] \& \&Y_{S'}\&\\
X\ \unsure{\alpha.|\text{\char38.}}\&\&\&\\Y\&\&\&
\end{tikzcd}}}
\[
 \widetilde w''(g\times_Yg')=\widetilde w(g)\times_Y\widetilde w'(g')
\]
moyennant l'identification
\[
 (X\times_YX')^{\{Y_{S'}\}}_Y\simeq X^{\{Y_{S'}\}}_Y\times_Y X'^{\{Y_{S'}\}}_Y
\]

\uline{Corollaire} Lois de composition compatibles avec des connexions.

% archive sheet number: 32

\scanpage{34}
Exemples,

Cas particuliers

% archive sheet number: 33

\scanpage{35}
(i)
\[
\begin{aligned}
 F&=\deleted{\unsure{X^p+Y^p-a}}aX^p+bY^p-1\\
 k[X,Y]/(F)&=A\\[-2pt]
 &\unsure{\wr\!\mid|\simeq}\\[-2pt]
 &k[Z]/(Z^p-a)[X]
\end{aligned}
\]
\margin{faire $Z=X+Y$ donc $F=Z^p-a$}
\margin{$aX^p+bY^p$}

\deleted{\unsure{transcendans}} $A$ est un anneau régulier, \inserted{non} $k$ est alg.\textsuperscript{t} fermé dans son corps des fractions

$A\otimes_k\bar k$ n'est pas intègre, il est isomorphe à $k[Z][T]/(T^p)$

\noindent\rule{.28\linewidth}{.3pt}

(ii)
\[
\left.
\begin{aligned}
 F&=Y^2-X^p+a\\
 k[X,Y]/(F)&=A
\end{aligned}
\right|\,p\ne2
\]
$A$ est un anneau régulier, \deleted{\unsure{$k$ est alg.\textsuperscript{t}}} son corps des fractions est extension séparable de $k$

$A\otimes_k\bar k$ n'est pas régulier, car le point $Y=0$, $X=a^{1/p}$ est singulier sur la courbe $F=0$

En effet, la courbe devient $F=Y^2-(X-\lambda)^p=0$ \quad$(\lambda=a^{1/p})$. Faisons $X-\lambda=Z$, on a donc $k[Y,Z]/(Y^2-Z^p)$. Or $Y^2-Z^p\in{\unsure{m|\mathfrak m}}^{2}$, où $m$ est l'idéal maximal de l'origine, donc on ne peut avoir point simple à l'origine. Cependant, pour $p\ne\unsure{2|1}$, $Y^2-Z^p$ est un polynôme irréductible. Cela prouve en $=$ temps que $K/k$ est une extension séparable. Montrons que $A$ est régulier.

\deleted{Dans la courbe obtenue par extension à $\bar k$, on voit que le seul point singulier est} \unsure{Car|le} .

$\Omega^1_k(A)\simeq A^2/A\bigl(\frac{\partial F}{\partial X},\frac{\partial F}{\partial Y}\bigr)\simeq A+A/2Y$. Les seuls points où $\Omega^1_k(A)$ soit non libre est aux idéaux premiers $\supset Y$, i.e. les idéaux premiers de $k[X,Y]$ \unsure{cont\textsuperscript{t}} $(Y,\deleted{\unsure{a}}Y^2-X^p+a)=(Y,X^p-a)$ ; or $k[X,Y]/(Y,X^p-a)\simeq k[X]/(X^{p}-a)$

Donc l'idéal $m'=(Y,X^p-a)$ est déjà maximal, c'est donc le

% archive sheet number: 34

\scanpage{36}
seul point singulier possible. Or, le carré de \inserted{${\unsure{m|\mathfrak m}}'/(F)$} est $m''/(F)$, \deleted{idéal \inserted{$m$ de $A$} \unsure{est}} où $m''=(Y^{\unsure{2}},Y(X^p-a),(X^p-a)^2,Y^2-X^p\unsure{-|+}a)$ \deleted{donc $Y=\frac{X^p}{\unsure{a}}\unsure{(}$} $=\bigl((X^p-a),Y^2,Y(X^p-a)\bigr)$, par suite $m/m^2=m'/m''$ est engendré par $Y$ (car $X^p-a$ est ds $m''$). Donc le point envisagé est aussi simple, mais non absolument simple.

\begin{center}\rule{.22\linewidth}{.3pt}\\[-.8ex]\rule{.19\linewidth}{.3pt}\end{center}

Dans le cas (i), il reste à vérifier que $k$ est algébriquement fermé dans $A$.

% archive sheet number: 35

\scanpage{37}
\margin{$\unsure{0|6}$\quad $\deleted{\Omega^1=0}$\par $A$}

$df=0$\qquad $f\in K^p$

\diagram{$\begin{array}{ll}
 \text{Critère de non ramification (Chevalley)} & \Omega^1_{B/A}=0\\
 \text{\rule{2em}{.3pt}\ d'étalement} & \Omega^1_{B/A}=0+\text{plat}\\[4pt]
 \text{Critère de \unsure{radicialité|radicalité}} & \Delta_{B/A}\ \text{est nilpotent}\\
 \text{\rule{2em}{.3pt}\ \unsure{distinct|distincts} (morph.\ \uline{fini})} &
 \Delta_{B/A}\ \text{est \uuline{nul}}
\end{array}$}

\margin{\uline{\unsure{Axiomatisation|Attention|Abréviation|Caractérisation}} les $\delta x_{\unsure{i|\cdot}}$ \unsure{engendrent} $\Delta$\par Cas \unsure{relatif|relient} : les $dx_i$ \unsure{engd} $\Omega$}

$\deltaf=0$ : $f$ est dans l'adhérence \deleted{différentielle} \deleted{\inserted{totale}} \deleted{\unsure{sur}} \unsure{universelle} de $A$ ds $B$

\deleted{$d\unsure{df=f|\delta f=\delta f}\ ?$} Si $B/A$ est plat : $f\in A$

Supposons que $B/A$ ait une $p$-base ; alors
\margin{$\Omega^1(B/A)$}
\[
\begin{aligned}
 df=0 &\Longleftrightarrow f\in A[B^p]\\
 \delta^{\unsure{(n)|(m)|(\infty)}}f=0 &\Longleftrightarrow f\in A[B^{p^{\unsure{k|n}}}]\\
 \delta^{(\infty)}f=0 &\Longleftrightarrow \delta f=0\Longleftrightarrow f\in A
\end{aligned}
\]

Cas d'extension \deleted{\unsure{séparables}} de corps \deleted{en \unsure{caract.} en car. $0$ :}

\diagram{\begin{tikzcd}[ampersand replacement=\&,column sep=small,row sep=tiny]
\&\& f\ \text{\unsure{alg. séparable}}\arrow[dl,Leftrightarrow] \\
df=0\arrow[rr,bend right=24,Rightarrow] \& \delta^{(\infty)}f=0\arrow[l,Rightarrow] \&[-1em] \deleted{\unsure{\Rightarrow}}\, f\ \text{algébrique}
\end{tikzcd}}
\margin{\unsure{si extension} \deleted{\unsure{alg.}} \deleted{\unsure{les}} en car. $0$}

% archive sheet number: 36

\scanpage{38}
\mspage{-9. 7-}
n° 215

opère à gauche. Désignons par $\Gamma_F$ le groupe de trivialité de $F$. Supposons que $F$ et $G/\Gamma_F$ soient munis de structures d'espèce $\unsure{\mathbb E|E}$ satisfaisant à la condition suivante :

C. \uline{si} $a\rightarrow s_a^*$ \uline{est une application dans} $G/\Gamma_{\handwritten{\unsure{F2|F1}}}$ \uline{d'un ensemble} $A$ \uline{muni d'une structure d'espèce} $\mathbb E$, \uline{une condition nécessaire et suffisante pour que cette application soit une} $\mathbb E$-\uline{application est que la famille d'applications} $(\xi\rightarrow s_a^*.\xi)_{a\in A}$ \uline{de} $F$ \uline{dans lui-même soit une} $\mathbb E$-\uline{famille}.

Il résulte en particulier de là que, pour tout $s^*\in G/\Gamma_F$, l'application $\xi\rightarrow s^*.\xi$ est un automorphisme de la structure d'espèce $\mathbb E$ de $F$. Nous allons montrer que la condition C. entraîne que l'espèce de structures $\unsure{\mathbb F|\mathbb E}(F;G)$ est plus puissante que $\mathbb E$. Soit en effet $F'$ un ensemble muni d'une structure d'espèce $\mathbb F(F;G)$, et soit $\Theta$ son ensemble de repères. Si $\theta\in\Theta$ et $s^*\in G/\Gamma_F$, nous désignerons par $\theta s^*$ le repère $\xi\rightarrow\theta(s^*.\xi)$. Tout élément $\theta$ de $\Theta$ permet de transporter à $F'$ la structure d'espèce $\mathbb E$ de $F$ ; tout autre élément de $\Theta$ étant de la forme $\theta s^*$, avec \handwritten{\inserted{un}} $s^*\in G/\Gamma_F$, la structure d'espèce $\mathbb E$ ainsi définie sur $F'$ ne dépend pas du choix de $\theta$. Les morphismes de l'espèce de structures $\mathbb F(F;G)$ étant des isomorphismes, il est clair que ce sont aussi des isomorphismes pour les structures d'espèce $\mathbb E$ de ces ensembles. On pourra donc parler d'espaces fibrés d'espèce $\mathbb F(F;G)$.

Soit maintenant $H$ un autre ensemble sur lequel $G$ opère à gauche et dont le groupe de trivialité $\Gamma_H$ contienne $\Gamma_F$ ; il y a alors une application canonique $s^*\rightarrow s^{**}$ de $G/\Gamma_F$ sur $G/\Gamma_H$. Supposons que $H$ satisfasse à la condition C. Ci-dessus (où on remplace $F$ par $H$), et que de plus la condition suivante soit satisfaite :

C'. \uline{l'application canonique de} $G/\Gamma_F$ \uline{dans} $G/\Gamma_H$ \uline{est une}

% archive sheet number: 37

\scanpage{39}
Foncteur \uline{différentiel} \uline{d'ordre $\unsure{n|r}$}

% archive sheet number: 38

\scanpage{40}
Foncteur différentiel d'ordre $\leq n$.

Soient $S$ un préschéma, $\unsure{C|\mathcal C}$ \deleted{\unsure{$X$|une}} une \inserted{s\textsuperscript{s}} catégorie fibrée pleine de la catégorie fibrée des préschémas simples de dim relative \unsure{$n$}, sur des $S'/S$, \inserted{les morphismes sur un $S'$ donné \unsure{étant} des \deleted{\unsure{isomorphismes}} \inserted{morphismes étales}} le plus souvent, $C$ est \unsure{tte} la catégorie en question.

Soit $\Phi$ une autre catégorie fibrée sur la $=$ base $\mathrm{Sch}_{/S}$ des préschémas sur $S$. [On pourra supposer dans les deux cas \deleted{\unsure{qu'ils}} une sous-catégorie de $(\mathrm{Sch})_{/\deleted{\unsure{S'}}S}$. \deleted{Peut-être faut-il plutôt prendre comme base $C$} \inserted{Elle définit une catégorie fibrée $\Phi^C$ sur $C$ par ``restriction''}]

Soit $F:C\longrightarrow\Phi_{\deleted{\unsure{C}}}^{C}$ une \deleted{foncteur fibrant} \inserted{section cartésienne} qui associe à tt $X'/S'\in C$ un $\xi_{X'/S'}\in\Phi_{\unsure{/C|/C'|\mathcal C}}$, de façon compatible \deleted{\unsure{avec}} avec les extensions de la base.
\margin{\diagram{\begin{tikzcd}[ampersand replacement=\&,column sep=small,row sep=small]
 \&\xi_{X'/S'}\&\xi_{X''/S''}\\
 \& X'\arrow[d,dash] \& X''\arrow[l]\arrow[d,dash]\\
 S \& S'\arrow[l] \& S''\arrow[l]
\end{tikzcd}}}

Ex : Le faisceau $\underline P^n_{X'/S'}$ [$\Phi$ la catégorie des Algèbres augmentées loc isom. \unsure{à} $\mathcal O[t_1,\ldots t_d]/(t)^{\unsure{n+1}}$] ou le fibré principal associé. Soit $\unsure{\delta_n|\mathscr D_n}$

Lorsqu'on se donne une catégorie fibrée une $\Psi$ sur \deleted{\unsure{$C$}} \inserted{$(\mathrm{Sch})/S$ \unsure{dont on tire|dont une} sur $C$ par restriction}, un élément $\xi\in\Psi_S$ et qu'on \deleted{\unsure{laisse de façon} $\operatorname{Aut}\bigl(\unsure{[\mathcal O_S]}[t_1,\ldots t_d]/(t_1,\ldots t_d)^{\unsure{n}}\bigr)\longrightarrow$} fait opérer \deleted{$\operatorname{Aut}\,\mathcal O_S[t_1,\ldots t_d]/(t_1,\ldots t_d)$} \inserted{le $G^{d,n}_S=\underline{\operatorname{Aut}}_S\bigl(\underline{\mathcal O}_S[t_1,\ldots t_d]/\unsure{(t_1,\ldots t_d)^{n+1}}$} sur $\xi$ et ses images inverses, [i.e. \deleted{\unsure{on}} a un

% archive sheet number: 39

\scanpage{41}
hom de \deleted{\unsure{fais}} foncteurs en groupes sur $(\mathrm{Sch})_{/S}$] on peut associer à $\xi$ un foncteur $F_\xi$, en utilisant $\unsure{\delta_n|\mathscr D_n}$.

On veut caractériser les foncteurs obtenus de cette façon. Notons que $\xi$ se \deleted{\unsure{trouve. vice}} \inserted{retrouve} à partir de $F_\xi$ en \deleted{\unsure{lui faisant}} \deleted{\inserted{faisant \unsure{opé}}} \inserted{\unsure{intégrant|interprétant}} \deleted{$S' = S,\ X\quad X'=E_S^n$ sur $S'$,} [\unsure{au dessus|en dessous}

\diagram{\begin{tikzcd}[ampersand replacement=\&,column sep=small,row sep=tiny,/tikz/execute at end picture={\draw[thin] (current bounding box.north west) -- (current bounding box.south east);}]
 \deleted{\xi_{E^n}} \& \deleted{\xi_{E^n_{S^{d,n}}}} \& \deleted{\xi} \& \\
 \& \deleted{\unsure{S'}\quad S_n=} \& \& \deleted{X=E_S^n}\\
 \deleted{E^n}\arrow[d] \& \deleted{E^n_{S^{d,n}}}\arrow[l]\arrow[d] \& \deleted{E^n}\arrow[l]\arrow[d] \& \deleted{S}\arrow[l]\\
 \deleted{S} \& \deleted{S^{d,n}}\arrow[l]\arrow[ul,dashed] \& \deleted{S}\arrow[l] \&\\
 \& \deleted{\Vert}\&\&\\
 \& \deleted{\operatorname{Spec}\mathcal O_S[t_1,\ldots t_d]/(t_1,\ldots t_d)^{n+1}}\&\&
\end{tikzcd}}

\deleted{de $S$ pour simplifier, mais \unsure{on répète|en le faisant} sur tt \unsure{$S'$|$T'$}]} \deleted{$G^{d,\unsure{n}}_{\deleted{\unsure{S}}}(S)$ comme $\underline{\operatorname{Aut}}_S(S^{d,n})$, que}
\inserted{\deleted{[augm $\varepsilon:S\longrightarrow S^{d,n}$]}}
\deleted{on \unsure{donne|laisse} opérer sur $E^n_{S^{d,n}}$ en laissant fixe $E^n$ considéré comme $E^n_{S^{d,n}}\times_{S^{d,n}}S$. Un tel automorphisme opère donc sur $\xi_{E^n_{S^{d,n}}}=\xi_{E^n}\times_S S^{d,n}$, de façon compatible avec les opérations sur $E^n_{S^{d,n}}$,}

% archive sheet number: 40

\scanpage{44}
LÊ Dũng Tráng\\
24 B\textsuperscript{d} Jules Ferry\\
75-Paris $\overline{\underline{\mathrm{XI}}}$\\

Paris le 11/01/68\\

à Monsieur le Professeur A. Grothendieck\\
IHES \qquad Bures S/Yvette\\

Monsieur le Professeur,\\

Ayant égaré les feuilles sur lesquelles étaient relevés les errata, aussi je me permets de vous envoyer si tardivement ces nouvelles-ci. Les errata ont été relevés pour la plupart dans les définitions. Je ne sais pas s'ils entravent des démonstrations car celles-ci faisaient intervenir des éléments de mathématiques que je \inserted{ne} possède pas comme les idéaux associés ou la profondeur d'un anneau. Pardonnez donc la trivialité de ces errata.

Veuillez agréer, Monsieur le professeur, l'expression\\
de mes sentiments respectueux et de ma profonde admiration.\\

\unsure{Tráng|Tg}.\\

% archive sheet number: 43

\scanpage{45}
\mspage{\textcircled{1}}
(20.1.3) \inserted{A la place de ``Nous allons nous intéresser... n'est autre que l'ensemble des éléments réguliers de $\mathcal O_{X,x}$'' lire :} ``Nous allons nous intéresser ici au cas où $\mathscr S$ est le sous-faisceau $\mathscr S(\mathcal O_X)$ de $\mathcal O_X$ tel que pour tout ouvert $U$, $\Gamma(U,\mathscr S)$ soit l'ensemble des éléments réguliers de l'anneau $\Gamma(U,\mathcal O_X)$ dont la restriction à tout ouvert $V\subset U$ est encore régulière dans $\Gamma(V,\mathcal O_X)$. Il est immédiat qu'il s'agit d'un faisceau car un élément est dans $\Gamma(U,\mathscr S)$ si et seulement si il est dans $\Gamma(U,\mathcal O_X)$ et \deleted{tous} ses germes en tout point $x$ de $U$ sont des éléments réguliers de $\mathcal O_{X,x}$. Il faut remarquer que $\mathscr S(\mathcal O_X)_x$ n'est pas \deleted{toujo} nécessairement l'ensemble des éléments réguliers de $\mathcal O_{X,x}$.''

(20.1.4) au lieu de : ``qui n'est autre que \deleted{un}\inserted{l'}anneau total des fractions de $\Gamma(U,\mathcal O_X)$'', -- lire : ``qui est un anneau de fractions à dénominateurs dans une partie multiplicativement stable composée d'éléments réguliers''.

(20.1.8) à la place de : ``en effet, si $s\in\Gamma(U,\mathscr S(\unsure{\mathcal M|\mathcal K}_X))$... inversible dans cet anneau de fractions'', -- lire : ``en effet, si $s\in\Gamma(U,\mathscr S(\mathcal M_X))$, pour tout $x\in U$ il existe un voisinage ouvert $V\subset U$ de $x$ tel que $s|V$ soit un élément régulier de $\Gamma(\deleted{\unsure{U}}V,\mathcal O_X)[\Gamma(V,\mathscr S)^{-1}]$ dont toutes les restrictions à des ouverts $W$ inclus dans $V$ sont des éléments réguliers de $\Gamma(W,\mathcal O_X)[\Gamma(W,\mathscr S)^{-1}]$ et on sait que dans ce cas $s|V$ est inversible dans $\Gamma(V,\mathcal O_X)[\Gamma(V,\mathscr S)^{-1}]$.''

(20.1.11) Remplacer ``$\mathscr S_f(U)$ l'ensemble des sections régulières $s\in\Gamma(U,\mathcal O_X)$ telles \inserted{que}...'', -- par : ``$\mathscr S_f(U)$ l'ensemble des sections régulières $s\in\Gamma(U,\mathscr S)$ \deleted{telles que} de $\mathcal O_X$ au dessus de $U$ telles que...''

% archive sheet number: 44

\scanpage{46}
\mspage{\textcircled{2}}
(21.1.5) A\deleted{u} la place de : ``Le faisceau $\mathscr S(\mathcal O_X)$ .... les éléments réguliers de $\Gamma(U,\mathcal O_X)$'', lire : ``Le faisceau $\mathscr S(\mathcal O_X)$ (20.1.3) dont les sections au-dessus d'un ouvert $U$ de $X$ sont les éléments réguliers de $\Gamma(U,\mathcal O_X)$ dont toutes les restrictions aux ouverts $V$ inclus dans $U$ sont encore des éléments réguliers''.

(21.1.6) A la fin du paragraphe au lieu de : ``il existe un élément régulier $t$ de $\Gamma(U,\mathcal O_X)$...'' lire ``il existe un élément \deleted{régulier} $t$ de $\Gamma(U,\mathscr S(\mathcal O_X))$''.

(21.2.9) Après la ligne (21.2.9.1) lire :

``qui est un isomorphisme : en effet si $U$ est un ouvert de $X$ tel que $\mathcal J|U=\mathcal O_U.f$ où $f\in\Gamma(U,{\unsure{\mathcal M|\mathcal K}}_X^*)$, d'après (21.2.2) on déduit que l'on a \deleted{un isomorphisme} restreint à $U$ un isomorphisme \deleted{de} qui à toute section $t$ de $\mathcal M_X(\mathcal J)|U$ au-dessus de $V\subset U$ fait correspondre une section $\bar t$ de $\mathcal M_X|U$ au-dessus de $V\subset U$ ; restreint à $U$ l'isomorphisme (21.2.9) est celui qui à toute section $t$ de $\mathcal M_X(\mathcal J)|U$ au-dessus de $V\subset U$ fait correspondre la section $\bar t(f|V)$ au-dessus de $V\subset U$. Dans l'isomorphisme (21.2.9,1)...''

(21.2.9) après la ligne (21.2.9.2) lire :

``et on désigne par $s_D$ la section méromorphe régulière de $\mathcal O_X(D)$ au-dessus de $X$ qui correspond par cet isomorphisme à la section $1$ de $\mathcal M_X$. Si $U$ est un ouvert de $X$ tel que $\mathcal O_X(D)|U=\mathcal O_U f^{-1}$ où $f\in\Gamma(U,\mathcal M_X^*)$ \deleted{et si} on a $D|U=\operatorname{div}_U(f)$. De plus pour un tel $U$ on a un isomorphisme de $\mathcal M_X(\mathcal J)|U$ sur $\mathcal M_X|U$ qui à toute section $t$ au-dessus de $V\subset U$ fait correspondre $\bar t$ ; cet isomorphisme fait correspondre à $s_D|U$ la section $f$ de $\Gamma(U,\mathcal M_X^*)$ donc :
\[
\operatorname{div}(s_D)=D.\tag*{(21,2.9.3)}
\]

% archive sheet number: 45

\scanpage{47}
\mspage{\textcircled{3}}

\deleted{(21,2.12) à la première ligne ajouter \deleted{\unsure{la|l'} hypothèse condition} : ``localement nœthérien'' entre ``préschéma'' et ``$X$''.}

\deleted{(21.3.6)}

(21.6.2) fin du paragraphe lire ${\mathfrak{Z}}^{p}(X)$ au lieu de $Z^{p}(X)$.

% archive sheet number: 46

\scanpage{49}

EGA $\overline{\underline{\mathrm{IV}}}$ 18

\diagram{\begin{tikzcd}[ampersand replacement=\&]
X \arrow[d,"\unsure{u|U}"']\\
Y \& \text{non ramifié}
\end{tikzcd}}

$U$ \deleted{\illegible{1}} ouvert schématiquement dense de $Y$

$Y_0$ sous-schéma fermé de $Y$, tel que $|Y_0|=|Y|$

$u_U$ une \deleted{\unsure{une}} section de $X$ sur $U$

$u_{Y_0}$ \rule[.5ex]{4.3cm}{.3pt} $Y_0$

$u_U$ et $u_{Y_0}$ coïncident sur $U_0=Y_0\cap U$.

Alors $\exists!$ section $u$ de $X$ sur $Y$, dont les restr. : $Y_0$ et $U$ sont $u_{Y_0}$ et $u_U$.

\noindent\rule{.52\linewidth}{.3pt}

\margin{EGA $\overline{\underline{\mathrm{VI}}}$}

\diagram{\begin{tikzcd}[ampersand replacement=\&]
X \arrow[d,"\text{non ram.}"] \& \& \\
{\underset{\text{\unsure{loc noeth|loc nœth|loc. noeth.}}}{Y}} \&
Y' \arrow[l,"\substack{\text{fini}\\\text{épim.}}"'] \&
Y'' \arrow[l,shift left=.45ex] \arrow[l,shift right=.45ex]
\end{tikzcd}}
$Y''\to Y'\mathbin{\times_Y}Y'$ surjectif

\margin{On doit pouvoir éliminer l'hyp. loc. noeth. et remplacer fini par \uline{entier}}

\diagram{\begin{tikzcd}[ampersand replacement=\&,column sep=small]
\text{Alors} \& X(Y)\arrow[r,dash] \& X(Y')\arrow[r,dash,shift left=.45ex]\arrow[r,dash,shift right=.45ex] \& X(Y'') \& \text{est exact}
\end{tikzcd}}

\noindent\rule{.7\linewidth}{.3pt}

\uline{Corollaire des deux} Variante en remplaçant $Y_0$ par un $Y'$ fini surjectif sur $Y$

\uline{Corollaire} Si $f:X\to Y$ est \unsure{ens.|ess.|un.} propre .....

% archive sheet number: 48

\scanpage{50}

\uline{Suites régulières}

Th. \unsure{II|11}.2.9. de Raynaud

% archive sheet number: 49

